% !TEX root = ../main.tex

\chapter{前言}

3D Gaussian Splatting 发展得太快了。原始 3DGS 把辐射场从慢速体渲染推向显式高斯与实时 rasterization 之后，后续工作很快不再只是“把图像渲染得更好”这一条线。它开始进入几何、材质、动态场景、4D、人像、机器人、自动驾驶、医学、遥感、安全隐私和生成式世界建模。论文数量增加之后，真正困难的事情反而变成了判断：哪些方法只是把 3DGS 放进一个新任务，哪些方法改动了关键模块，哪些方法会成为后续工作的局部锚点。

这本书就是为这个判断而写的。它不是原有长版技术综述的替代品，也不是短版的再压缩版。长版更像完整地图，尽可能覆盖 3DGS 研究生态；短版帮助读者快速建立方向感；而这本“可信方法综述”只关注那些至少具备一项可核验外部信息的方法，例如知名企业参与、公开代码、GitHub 上已有较高关注度、或者论文引用已经明显累积。这样做不是为了给论文打分，而是为了在论文过多、锚点不可靠的情况下，先把读者带到更值得细读、复现和持续跟踪的方法上。

本书的写法沿用原技术综述中最有用的部分：从问题出发，而不是从论文列表出发。每个方向先说明原始 3DGS 或上一代锚点方法卡在哪里，再说明新的技术路线如何改动表示、初始化、密度控制、渲染、优化、损失、压缩、动态建模、生成流程或应用接口。锚点方法会写得更厚，因为读者需要知道它为什么能成为后续方法的参照；其他方法则围绕锚点说明差异，避免重复讲同一套背景。

这本书也尽量把“方法到底怎么做”说清楚。只说“提升质量”“增强鲁棒性”“改善效率”是不够的。一个方法如果改变了高斯的参数化，就要说明参数接在哪里；如果改了渲染器，就要说明排序、混合或可微传播如何变；如果引入扩散模型、视频模型或语言模型，就要说明它们提供的是初始化、先验、监督、补全还是编辑接口；如果是系统或应用论文，也要说明 3DGS 在整个系统中承担的是地图、记忆、渲染、规划接口还是可优化中间表示。

书中随方法给出企业背景、代码链接、GitHub 关注度和引用情况。写明这些信息，是为了让读者知道一个方法是否已经有工程复现条件、是否受到社区关注、是否有企业或研究机构投入，而不是把它们当成绝对的学术价值判断。很多新论文很有路线意义，但代码还没放出，引用也还没来得及积累；这类工作如果已经能说明技术路线，本书仍会收录，只会在正文中诚实说明当前可核验信息。

\enlargethispage{4\baselineskip}
最后，这本书应当被看成一个持续更新的技术档案。3DGS 生态仍在高速变化，新论文会继续出现，旧论文的代码会补齐，引用和星标也会变化。后续更新不应该简单追加论文摘要，而应该继续坚持本书的增量写法：先找锚点，再讲差异；先讲问题，再讲机制；先确认可核验信息，再把方法放入合适的技术路线。

\section*{2026 年 11 月增补范围}

本次以论文库原有入选条件为起点，补核了 108 条此前未写入正文的记录。其中 101 篇补入各章，7 篇因唯一外部线索无法证实而留在代码待跟进表。本版统计正文方法标记，一篇论文只计一次。外部线索只说明入选理由，实验表现仍以具体任务、数据和对比条件为准。

\chapter{如何阅读本书}

\section*{先看锚点，再看增量}

如果你刚进入 3DGS 领域，建议先读第 1 章，理解原始 3DGS 的显式高斯表示、可微 rasterization 和自适应密度控制。之后不要急着按时间线追论文，而是按问题读：想理解几何和表面，就读 2DGS、几何正则与表面重建相关章节；想理解动态，就读 4D 和单目动态路线；想做系统落地，就读压缩、大规模、SLAM、机器人和自动驾驶章节。

每章开头都有“本章阅读锚点”。它的作用是告诉你这一章主要围绕哪些基础环节展开：表示、渲染、优化、密度控制、几何、动态、压缩、生成或应用接口。读具体方法时，可以先问一句：它是在替换哪个环节，还是在把 3DGS 接入一个新的任务系统？

\section*{如何看方法段落}

本书的每个方法段落都尽量回答四个问题。第一，它继承哪条路线，是原始 3DGS、Scaffold-GS、Mip-Splatting、2DGS、4DGS、GaussianAvatar，还是某个更局部的锚点。第二，原方法在什么条件下不够好，是几何不准、动态不稳、训练太慢、内存太大、泛化差，还是应用接口不完整。第三，它具体改了什么模块。第四，它怎样训练、渲染或推理。

方法信息中的企业、代码、GitHub stars 和引用信息只作为阅读线索。看到代码链接，说明当前材料中已经能核验到可用仓库；看到“星标高”或“引用很多”，说明它在社区或文献中已经形成了较明显的外部关注。反过来，没有这些信息并不等于方法不重要，尤其是 2026 年的新论文，很多还处在代码和引用尚未稳定的阶段。

\section*{如何使用待跟进清单}

附录中的代码待跟进表列出了一批已经出现公开链接、但当前数据库还不能确认完整训练或推理代码的工作。它们没有被当作“已有完整代码”写进正文方法信息，但值得后续复核。对工程读者来说，这张表可以作为持续跟踪清单；对综述更新来说，它是下一轮外部信息更新流程的优先检查对象。
